\part{The Interface}

\chapter{Screen Layout}

FairWindSK keeps three shell areas visible at all times.
These names are used consistently throughout this manual, in the project source code,
and in all contributor documentation.

\section{The Three Zones}

\begin{center}
\begin{tcolorbox}[enhanced,colback=LtBlue,colframe=Navy,boxrule=1.5pt,arc=4pt,
  top=7pt,bottom=5pt,left=12pt,right=12pt,width=0.96\textwidth]
  {\sffamily\bfseries\color{Navy}TOP BAR --- always visible}\\[3pt]
  {\small\icon{lcd-app}~App icon \quad
         \icon{layout-signalk}~Data widgets \quad
         App name \quad
         \icon{lcd-clock}~Date/time \quad
         Status icons}
\end{tcolorbox}
\vspace{2pt}
\begin{tcolorbox}[enhanced,colback=white,colframe=Navy,boxrule=1.5pt,arc=4pt,
  top=26pt,bottom=26pt,left=12pt,right=12pt,width=0.96\textwidth]
  \centering
  {\large\sffamily\bfseries\color{MidGray}APPLICATION AREA}\\[4pt]
  {\small\sffamily\color{MidGray}\itshape
    Launcher \textbar{} Chart plotter \textbar{} Instruments \textbar{}
    Settings \textbar{} My~Data \textbar{} Signal~K web apps}
\end{tcolorbox}
\vspace{2pt}
\begin{tcolorbox}[enhanced,colback=LtBlue,colframe=Navy,boxrule=1.5pt,arc=4pt,
  top=7pt,bottom=5pt,left=12pt,right=12pt,width=0.96\textwidth]
  {\sffamily\bfseries\color{Navy}BOTTOM BAR --- always visible}\\[3pt]
  {\small
    Open web apps strip \quad
    \icon{applications}~Apps \quad
    \icon{alarm-pob}~POB \quad
    \icon{command-autopilot}~Autopilot \quad
    \icon{anchor-iec}~Anchor \quad
    \icon{alarm-general}~Alarms \quad
    \icon{database}~My~Data \quad
    \icon{settings-iec}~Settings \quad
    \icon{layout-signalk}~Status}
\end{tcolorbox}
\end{center}

\screenshot{screen-layout-overview}{Screen layout: Top Bar, Application Area, Bottom Bar}{}

\section{The Application Area}

The Application Area is the main working space between the bars.
This is where FairWindSK shows:

\begin{itemize}
  \item launcher pages
  \item chart and navigation apps
  \item Settings
  \item My~Data
  \item embedded Signal~K web applications
\end{itemize}

The active view must not visually overlap the Top Bar or Bottom Bar unless a
formally defined drawer is temporarily covering it.

\section{Drawers}

\rowcolors{2}{LtGray}{white}
\begin{longtable}{L{4.0cm}L{9.7cm}}
\toprule
\rowcolor{Navy}
\color{white}\sffamily\bfseries Drawer &
\color{white}\sffamily\bfseries Description \\
\midrule
Bottom Bar horizontal drawer &
  Modal dialog opening upward from the Bottom Bar.
  Used for web modal dialogs, confirmations, file pickers, and settings sub-panels.
  Only one dialog is visible at a time; additional dialogs stack.
  Can extend upward to temporarily cover the Application Area.
  The hosted UI must fit inside the drawer without requiring vertical scrolling. \\
Application Area vertical drawer &
  Non-modal panel opening leftward from the right edge of the Application Area.
  Used by apps for contextual menus or tool palettes.
  The Application Area remains interactive behind it. \\
\bottomrule
\end{longtable}

\section{Single-Window Model}

FairWindSK enforces a single-window design on all platforms.
No app opens a second window; no external native application is launched.
\code{file://} launcher entries are blocked.

Android is the intentional platform exception for selected native applications.
On Android 13/API 33 or newer, FairWindSK can optionally act as the Home app and
launch activities explicitly selected in \ui{Settings~> Android}. Those activities
temporarily leave the FairWindSK window and return through Android Back or Home;
all FairWindSK and Signal~K content continues to use the single-window shell.

On desktop platforms (Linux, macOS, Windows, Raspberry~Pi~OS),
press \key{SHIFT+TAB} to bring FairWindSK back to the foreground if a hosted web
application has taken full focus.

\chapter{The Top Bar}
\label{ch:topbar}

The Top Bar is always visible.
Its purpose: provide a two-second confidence check at any moment.
Use it to answer: \textit{What is happening right now, and is the data still
trustworthy?}

\screenshot{topbar-overview}{Top Bar: position, SOG, DPT, BTW, ETA and status icons}{}

\section{Fixed Top Bar Elements}

\rowcolors{2}{LtGray}{white}
\begin{longtable}{C{1.5cm}L{2.8cm}L{9.4cm}}
\toprule
\rowcolor{Navy}
\color{white}\sffamily\bfseries Icon &
\color{white}\sffamily\bfseries Element &
\color{white}\sffamily\bfseries Function \\
\midrule
\iconlg{lcd-app}       & App icon    &
  Icon of the current application.
  Tapping opens the Application Area vertical drawer. \\
---                     & App name    &
  Name of the currently active application. \\
\iconlg{lcd-clock}     & Date / time &
  Current device time and date. \\
\iconlg{layout-signalk}& Status icons &
  REST and stream health indicators. \\
\bottomrule
\end{longtable}

\section{Data Widgets --- Complete Reference}

Each widget subscribes to a Signal~K path configured in
\ui{Settings~> Signal~K Paths}.

\rowcolors{2}{LtGray}{white}
\begin{longtable}{C{1.4cm}C{1.1cm}L{4.4cm}L{3.0cm}L{3.8cm}}
\toprule
\rowcolor{Navy}
\color{white}\sffamily\bfseries Icon &
\color{white}\sffamily\bfseries Key &
\color{white}\sffamily\bfseries Default Signal~K path &
\color{white}\sffamily\bfseries Displays &
\color{white}\sffamily\bfseries Operational significance \\
\midrule
\iconlg{lcd-position} & \code{pos} & \skpath{navigation.position}                              & Latitude / longitude      & GPS fix confirmation \\
\iconlg{lcd-cog}      & \code{cog} & \skpath{navigation.courseOverGroundTrue}                  & COG (\textdegree{}T)      & Actual track over seabed \\
\iconlg{lcd-sog}      & \code{sog} & \skpath{navigation.speedOverGround}                       & SOG (kn)                  & True speed for ETAs \\
\iconlg{lcd-hdg}      & \code{hdg} & \skpath{navigation.headingTrue}                           & Heading (\textdegree{}T)  & Where the bow points \\
\iconlg{lcd-stw}      & \code{stw} & \skpath{navigation.speedThroughWater}                     & STW (kn)                  & Sail trim, current deduction \\
\iconlg{lcd-dpt}      & \code{dpt} & \skpath{environment.depth.belowTransducer}                & Depth below transducer    & Safety-critical in shoal water \\
\iconlg{lcd-btw}      & \code{btw} & \skpath{navigation.course.calcValues.bearingTrue}         & Bearing to WPT            & Course to steer \\
\iconlg{lcd-dtg}      & \code{dtg} & \skpath{navigation.course.calcValues.distance}            & Distance to go            & Passage progress \\
\iconlg{lcd-ttg}      & \code{ttg} & \skpath{navigation.course.calcValues.timeToGo}            & Time to go                & ETA planning \\
\iconlg{lcd-eta}      & \code{eta} & \skpath{navigation.course.calcValues.estimatedTimeOfArrival} & Estimated arrival      & Crew notification \\
\iconlg{lcd-xte}      & \code{xte} & \skpath{navigation.course.calcValues.crossTrackError}     & Cross-track error         & Route compliance \\
\iconlg{lcd-vmg}      & \code{vmg} & \skpath{performance.velocityMadeGood}                     & Velocity made good        & True sailing efficiency \\
\iconlg{lcd-wpt}      & \code{wpt} & \skpath{navigation.course.nextPoint}                      & Active waypoint name       & Route awareness \\
\bottomrule
\end{longtable}

\begin{warningbox}
Missing depth is \textbf{not} zero depth.
Missing means the measurement is unavailable --- treat it as entirely unknown.
If depth becomes stale or missing in shallow or restricted water, verify
immediately and adjust speed and manoeuvring margin accordingly.
\end{warningbox}

\section{Reading Navigation Values Correctly}

\subsection{COG and SOG}

COG is course over ground; SOG is speed over ground.
These show how the boat is actually moving over the earth, which may differ from
where the bow is pointing due to leeway and tidal current.

\subsection{Heading}

Heading tells you where the bow is pointed.
If heading becomes stale or missing, steer conservatively and cross-check with a
trusted compass.

\subsection{Route and waypoint values}

When routing is active, BTW, DTG, TTG, ETA, XTE, and VMG all update continuously.
If any of these values stop updating, check the chart and your surroundings before
relying on the route guidance.

\section{Configuring the Top Bar}

Open \ui{Settings~> Top Bar}.
A live preview of the bar appears at the top; the widget palette is below.
Use the \iconlg{arrow-left-google}~move-left and \iconlg{arrow-right-google}~move-right
buttons to reposition widgets.
Toggle \key{Show Icon}, \key{Show Text}, \key{Show Units}, and \key{Show Trend}
per widget.
Tap \iconlg{refresh-google}~\key{Reset Defaults} to restore the built-in layout.

\begin{tipbox}
In rough conditions, a lean Top Bar with 3--4 key widgets (position, SOG, DPT, BTW)
is safer than twelve items.
Add widgets only if they add genuine situational awareness at the current passage
phase.
\end{tipbox}

\chapter{The Bottom Bar}
\label{ch:bottombar}

The Bottom Bar is the permanent action strip at the bottom of every screen.
Its spatial stability is intentional: the \iconlg{applications}~Apps button is always
in the same place, the \iconlg{alarm-pob}~POB button is always within reach.
Muscle memory is safety.

\screenshot{bottombar-overview}{Bottom Bar: open-apps strip, main icons, Signal~K status}{}

\section{Bottom Bar Regions}

\rowcolors{2}{LtGray}{white}
\begin{longtable}{L{3.6cm}L{10.1cm}}
\toprule
\rowcolor{Navy}
\color{white}\sffamily\bfseries Region &
\color{white}\sffamily\bfseries Description \\
\midrule
Open web-apps strip &
  Scrollable row of small app icons on the left, one per loaded application.
  Tap any icon to jump to that application instantly.
  The active application is highlighted. \\
Main icons &
  The central group of seven fixed icons.
  Spatially stable across all screens and comfort presets. \\
Signal~K status area &
  Right side. Compact stream and REST health indicators showing server name,
  last update timestamp, and connection state.
  Tapping opens a status detail drawer. \\
\bottomrule
\end{longtable}

\section{Main Icon Reference}

These are the exact icons defined in \code{BottomBar.ui}.

\rowcolors{2}{LtGray}{white}
\begin{longtable}{C{1.4cm}L{2.4cm}L{5.4cm}L{4.5cm}}
\toprule
\rowcolor{Navy}
\color{white}\sffamily\bfseries Icon &
\color{white}\sffamily\bfseries Label &
\color{white}\sffamily\bfseries Function &
\color{white}\sffamily\bfseries Notes \\
\midrule
\iconlg{applications}     & Apps      &
  Opens the launcher. Hides the current application and shows the tile grid. &
  Use as the safe escape hatch from any stuck or confusing screen. \\
\iconlg{alarm-pob}        & POB       &
  Person Over Board. One tap starts the emergency. &
  Always visible. Never disabled. \\
\iconlg{command-autopilot}& Autopilot &
  Opens the Autopilot bar. &
  Grey/inactive if no compatible autopilot plugin is detected. \\
\iconlg{anchor-iec}       & Anchor    &
  Opens the Anchor Watch bar. &
  Grey/inactive if no anchor plugin is detected. \\
\iconlg{alarm-general}    & Alarms    &
  Opens the Alarms bar. Highlighted when any alarm is active. &
  Accent colour when active; draws the eye across all presets. \\
\iconlg{database}         & My~Data   &
  Opens the resource manager (Chapter~\ref{ch:mydata}). & \\
\iconlg{settings-iec}     & Settings  &
  Opens all configuration (Chapter~\ref{ch:settings}). & \\
\bottomrule
\end{longtable}

\section{Operational Overlay Bars}

Tapping \iconlg{alarm-pob}~POB, \iconlg{command-autopilot}~Autopilot,
\iconlg{anchor-iec}~Anchor, or \iconlg{alarm-general}~Alarms replaces the normal Bottom
Bar icons with a dedicated operational bar.
The Application Area continues running behind the bar.
Tap \iconlg{close-google}~close (rightmost icon) to return to the normal Bottom Bar
without changing any operational state.

\begin{warningbox}
Operational bars are not modal.
Data continues updating.
The chart plotter and all applications remain live behind the operational bar.
\end{warningbox}

\chapter{The Launcher and Application Management}
\label{ch:launcher}

The launcher is the safe home position of the shell.
Use it to return to app tiles, switch between tasks, leave a frozen app, or
reorient yourself after a confusing workflow.
Whenever the current screen stops feeling clear or trustworthy, tap
\iconlg{applications}~\key{Apps} to return here.

\screenshot{launcher-overview}{Launcher home screen: tile grid, page indicator, Bottom Bar}{}

\section{Opening an Application}

\begin{enumerate}
  \item Tap \iconlabel{applications}{Apps} in the Bottom Bar.
  \item Navigate pages by swiping if needed.
  \item Tap the tile for the application you want.
  \item Wait for it to load in the Application Area.
\end{enumerate}

If an app does not appear correctly, avoid repeated rapid taps.
Wait briefly, retry once if the shell offers the option, then return to the launcher
if the app still looks wrong.

\section{Switching Between Open Applications}

Use the open-apps strip in the Bottom Bar to switch between loaded applications
with a single tap.
The strip shows a small icon for each loaded application; the active one is
highlighted.

\section{Adding a Custom Application}

\begin{enumerate}
  \item Open \ui{Settings~> Applications}.
  \item Tap \iconlabel{define-new-application-in-palette}{Create app}.
  \item Enter a display name and the full URL (\code{http://} or \code{https://}).
  \item Optionally set a custom icon path and zoom level.
  \item Tap \key{Save}.
\end{enumerate}

\section{Recovering a Stuck or Blank Application}

\begin{enumerate}
  \item Wait briefly. Retry once if the shell offers a retry action.
  \item Tap \iconlabel{applications}{Apps} to return to the launcher.
  \item Check shell health and connection state.
  \item Reopen the app only after confirming the shell itself is healthy.
\end{enumerate}

\begin{tipbox}
A failed application does not mean FairWindSK has failed.
The shell, data widgets, and operational bars (POB, anchor, autopilot) continue
working independently of any individual application.
A degraded app is not the same as a degraded vessel state.
\end{tipbox}
